\documentclass{article} %210 mm × 297 mm
\usepackage[utf8]{inputenc}
\usepackage[a4paper,left=1.5cm, right=1.5cm, top=2cm, bottom=2cm]{geometry}
\usepackage{graphicx}
%\usepackage{blindtext}
\usepackage{multirow}
\usepackage{mhchem}
\usepackage{url} 
\usepackage{subfigure}
\usepackage{amsfonts,latexsym} % para tener disponibilidad de diversos simbolos
\usepackage{enumerate}
%\usepackage{booktabs}
\usepackage{float}
%\usepackage{threeparttable}
\usepackage{array,colortbl}
%\usepackage{ifpdf}
%\usepackage{rotating}
\usepackage{cite}
\usepackage{stfloats}
\usepackage{url}
\usepackage{listings}
\usepackage{fancyhdr}

\usepackage{color}
%\usepackage{hyperref}
%\usepackage{wrapfig}
\usepackage{array}
%\usepackage{multirow}
%\usepackage{adjustbox}
%\usepackage{nccmath}
\usepackage[dvipsnames]{xcolor}


\newcommand{\titulo}{Abgabe 4; MfN II.}
\newcommand{\nombre}{Liliana Sanfilippo}
\newcommand{\FCE}{\textbf{SoSe 2025}}

 

\usepackage{fancyhdr}
\pagestyle{fancy}
\fancyhead{}
\fancyfoot{}
\fancyhead[LO]{\small\nombre}
\fancyhead[RE]{\small\titulo}
\fancyhead[LO]{\small\FCE}
\fancyfoot[RO,LE] {\thepage}

\setcounter{page}{1} %Número de la página, la editorial lo cambiará



\usepackage[onehalfspacing]{setspace}
\begin{document}


\begin{tabular}{p{12cm}}
{{\Large\bf\titulo}}\\ 
\vspace{0.2cm}\\
 {\large\nombre}\\
\end{tabular}
\vspace{0.2cm}\\
\section*{Aufgabe 4.1}
\subsection*{(a)}
\begin{align*}
      A & = \begin{pmatrix}
    -2 & 2 & -8 & -6 \\ 
    0 & 0 & -8 & -6 \\
    0 & 0 & 4 & 3 \\ 
    0 & 0 & -6 & -5
    \end{pmatrix} & \mid z1=z1-z2\\ 
    & =  \begin{pmatrix}
    -2 & 2 & 0 & 0 \\ 
    0 & 0 & -8 & -6 \\
    0 & 0 & 4 & 3 \\ 
    0 & 0 & -6 & -5
    \end{pmatrix} & \mid z2=z2 + 2\cdot z3\\ 
     & =  \begin{pmatrix}
    -2 & 2 & 0 & 0 \\ 
    0 & 0 & 0 & 0 \\
    0 & 0 & 4 & 3 \\ 
    0 & 0 & -6 & -5
    \end{pmatrix}\\ 
\end{align*}
\[
      A - \lambda I  = \begin{pmatrix}
    -2 - \lambda & 2 & 0 & 0\\ 
    0 &- \lambda & 0 & 0\\
    0 & 0 & 4 - \lambda  & 3 \\
    0 & 0 & -6 & -5 - \lambda
    \end{pmatrix}
\]
Laplace Entwicklung nach der ersten Spalte. 
\begin{align*}
 det( A - \lambda I) & = \sum^4_{j=1} a_{1j} \cdot (-1)^{1+j} \cdot det(A'_{1j}) \\
    & = a_{11} \cdot (-1)^{2} \cdot det(A'_{11})\\
    & =(-2-\lambda) \cdot det(A'_{11}) \\ \\
    A'_{11} & =  \begin{pmatrix}
    - \lambda & 0 & 0 \\
     0 & 4 - \lambda  & 3 \\ 
     0 & -6 & -5 - \lambda 
    \end{pmatrix} \\ \\
\end{align*}
Laplace nach der ersten Spalte. 
\begin{align*}
    det(A'_{11}) & =  \sum^3_{j=1} a_{1j} \cdot (-1)^{1+j} \cdot det(A''_{1j}) \\
    & = a_{11} \cdot (-1)^{2} \cdot det(A''_{11}) \\
     & = (-\lambda) \cdot det(A''_{11}) \\ \\
     \Rightarrow  det( A - \lambda I) &  =(-2-\lambda) \cdot  (-\lambda \cdot det(A''_{11}) )  \\ \\
     A''_{11} & =  \begin{pmatrix}
     4 - \lambda  & 3 \\
      -6 & -5 - \lambda
    \end{pmatrix} \\ \\
    det ( A''_{11}) & = (4-\lambda) \cdot (-5-\lambda) -  3 \cdot (-6) \\
    & = \lambda^2 + \lambda - 2 
\end{align*}
\begin{align*}
      \Rightarrow  det( A - \lambda I) &  =(-2-\lambda) \cdot  (-\lambda \cdot (\lambda^2 + \lambda - 2 ) )  \\ 
     &  =(-2-\lambda) \cdot  (-\lambda^3 - \lambda^2 + 2 \lambda ) \\
     & = \lambda^4 + 3\lambda^3 - 4 \lambda 
\end{align*}
Also: 
 \begin{align*}
 \chi_A(\lambda) &  = \lambda^4 + 3\lambda^3 - 4 \lambda  \\
 & = \lambda \cdot  (\lambda^3 + 3\lambda^2 - 4)  \\  \\ 
  0 &    = 0 \cdot  (0^3 + 3 \cdot 0^2 - 4) \\
  \Rightarrow & \lambda_1 = 0 \\ \\
  0 & = 1^3 + 3\cdot 1^2 - 4 \\ 
  \Rightarrow & \lambda_2 = 1 \\ \\
  \text{Zerlegen} \\ 
  \lambda^3 + 3\lambda^2 - 4  & = (\lambda - 1) (\lambda^2 + 4 \lambda + 4)  \\
  & = (\lambda - 1) (\lambda + 2) (\lambda + 2) \\ \\
  \Rightarrow & \lambda_3 = -2 \\
  & \lambda_4 = -2
 \end{align*}
Eigenvektoren und deren Basen berechnen.   
\paragraph{Für $\lambda_1$:}
\begin{align*}
        A - 0I &  = \begin{pmatrix}
    -2 - 0 & 2 & 0 & 0\\ 
    0 &- 0 & 0 & 0 \\
    0 & 0 & 4 - 0  & 3 \\
    0 & 0 & -6 & -5 - 0
    \end{pmatrix} \\
   &  = \begin{pmatrix}
    -2 & 2 & 0 & 0 \\ 
    0 &0 &0 & 0 \\
    0 & 0 & 4   & 3 \\
    0 & 0 & -6 & -5
    \end{pmatrix}
\end{align*}
    Setze $ (A - 0I)v = 0$.
\[\begin{pmatrix}
    -2 & 2 & 0 & 0 \\ 
     0 & 0 & 0 & 0\\
    0 & 0 & 4  & 3 \\
    0 & 0 & -6 & -5 
    \end{pmatrix} \begin{pmatrix}
       a \\ b \\ c \\ d 
    \end{pmatrix}  = \begin{pmatrix}
       0 \\ 0 \\0 \\ 0
    \end{pmatrix}\]
    \[  \begin{array}{cccccc}
       -2a  & +2b &  &  & = 0  \\
       &&  4c & +3d & = 0 \\
        && -6c & -5d & = 0
    \end{array} \] 
    \[  \Rightarrow a = b   \quad  und \  \Rightarrow 4c = -3d \quad  
   \Rightarrow c = -\frac{3}{4}d \]
\begin{align*}
   -6 \cdot (-\frac{3}{4})d -5d  & = 0 \\
    \Rightarrow  4,5d -5d & = 0 \\
   \Rightarrow c = d = 0 
\end{align*}
\begin{align*}
       \Rightarrow & \quad \text{ EV haben die Form } \begin{pmatrix}
         a \\ a \\ 0 \\ 0
      \end{pmatrix} \\
        \Rightarrow & \quad \text{Basis des Eigenraums für } \lambda = 0 : \Biggl\{ \begin{pmatrix}
         1 \\ 1 \\ 0 \\ 0
      \end{pmatrix}   \Biggr\} 
\end{align*}   
\paragraph{Für $\lambda_2$:}
\begin{align*}
       A - 1I &  = \begin{pmatrix}
    -3 & 2 &0 & 0\\ 
    0 &- 1 & 0 & 0 \\
    0 & 0 & 3  & 3 \\
    0 & 0 & -6 & -6
    \end{pmatrix} & \mid z1 = z1+ 2\cdot z2 \quad und  \quad z4 = z4 + 2\cdot z3\\
    & = \begin{pmatrix}
    -3 & 0 &0 & 0\\ 
    0 &- 1 & 0 & 0 \\
    0 & 0 & 3  & 3 \\
    0 & 0 & 0 & 0
    \end{pmatrix} 
\end{align*}
    Setze $ (A - 1I)v = 0$.
    \[   \begin{pmatrix}
    -3 & 0 & 0 &  0\\ 
    0 &- 1 & 0 & 0 \\
    0 & 0 & 3  & 3 \\
    0 & 0 & 0 & 0
    \end{pmatrix} \begin{pmatrix}
       a \\ b \\ c \\ d 
    \end{pmatrix}  = \begin{pmatrix}
       0 \\ 0 \\0 \\ 0
    \end{pmatrix} \]
\begin{align*}
    \begin{array}{cccccc}
       -3a  &  &  & & = 0  \\
       & -b &   &  & = 0 \\
       &&  3c & +3d & = 0 \\
    \end{array} \\
    %%%%%%%%%%%%%%%%%%
     \Rightarrow c & = -d \\
     \Rightarrow  a & = b = 0
\end{align*}
\begin{align*}
      \Rightarrow & \quad \text{ EV haben die Form } \begin{pmatrix}
         0 \\ 0 \\ -d \\ d
      \end{pmatrix}  \\
       \Rightarrow & \quad \text{Basis des Eigenraums für } \lambda = 1 : \Biggl\{ \begin{pmatrix}
          0 \\ 0 \\ -1 \\ 1
      \end{pmatrix}   \Biggr\}
\end{align*}
\paragraph{Für $\lambda_3$ und $\lambda_4$:}
\begin{align*}
     A +2I &  = \begin{pmatrix}
    0 & 2 & 0 & 0 \\ 
    0 & 2 & 0 & 0\\
    0 & 0 & 6  & 3 \\
    0 & 0 & -6 & -3
    \end{pmatrix} & \mid z2 = z2-z1 \quad und \quad z4 = z4+z3 \\ 
    &  = \begin{pmatrix}
    0 & 2 & 0 & 0 \\ 
    0 & 0 & 0 & 0\\
    0 & 0 & 6  & 3 \\
    0 & 0 & 0 & 0
    \end{pmatrix}
\end{align*}
    Setze $(A +2 I)v = 0$.
\begin{align*}
       \begin{pmatrix}
     0 & 2 & 0 & 0 \\ 
    0 & 0 & 0 & 0\\
    0 & 0 & 6  & 3 \\
    0 & 0 & 0 & 0
    \end{pmatrix} \begin{pmatrix}
       a \\ b \\ c \\ d 
    \end{pmatrix} & = \begin{pmatrix}
       0 \\ 0 \\0 \\ 0
    \end{pmatrix}  
     \end{align*}
     \begin{align*}
         \begin{array}{cccccc}
             & 2b &  &  & =0\\
    &  & 6c  & +3d & = 0 
         \end{array} \\
         \Rightarrow b = 0\\
         \Rightarrow d = -2c \\
          \Rightarrow a \text{ frei}
     \end{align*}
     \[
       \Rightarrow \quad \text{ EV haben die Form }  \begin{pmatrix}
         a \\ 0 \\ c \\ -2c 
      \end{pmatrix} \quad 
        also \quad  a \cdot  \begin{pmatrix}
         1 \\ 0 \\ 0 \\ 0 
      \end{pmatrix} + c \cdot  \begin{pmatrix}
         0 \\ 0 \\ 1 \\ -2 
      \end{pmatrix}
     \]
     \begin{align*}
       \Rightarrow  \quad \text{Basis des Eigenraums für } \lambda = -2 : & \Biggl\{ \begin{pmatrix}
         1 \\ 0 \\ 0 \\ 0 
      \end{pmatrix}, \begin{pmatrix}
         0 \\ 0 \\ 1 \\ -2 
      \end{pmatrix}   \Biggr\}
     \end{align*}
 \subsection*{(b)}
\[
B = ( \begin{pmatrix}
         1 \\ 1 \\ 0 \\ 0 
      \end{pmatrix}, \begin{pmatrix}
         0 \\ 0\\ -1 \\ 1
      \end{pmatrix}, \begin{pmatrix}
         1 \\ 0 \\ 0 \\ 0
      \end{pmatrix}  , \begin{pmatrix}
           0 \\ 0 \\ 1 \\ -2
      \end{pmatrix}  )
\]

\[
S =  \begin{pmatrix}
         1&0&1&0\\
         1&0&0&0 \\
         0&-1&0&1 \\
         0&1&0&-2
      \end{pmatrix}
\]

\subsection*{(c)}
\[
A \cdot S = \begin{pmatrix}
    -2 & 2 & 0 & 0 \\ 
    0 & 0 & 0 & 0 \\
    0 & 0 & 4 & 3 \\ 
    0 & 0 & -6 & -5
    \end{pmatrix} \cdot   \begin{pmatrix}
        1&0&1&0\\
         1&0&0&0 \\
         0&-1&0&1 \\
         0&1&0&-2
      \end{pmatrix} = 
      \begin{pmatrix}
   0&0&-2&0\\
           0&0&0&0\\
           0&-1&0&-2\\
           0&1&0&4
      \end{pmatrix}
\]
\[
S \cdot D =  \begin{pmatrix}
       1&0&1&0\\
         1&0&0&0 \\
         0&-1&0&1 \\
         0&1&0&-2
      \end{pmatrix} \cdot  
      \begin{pmatrix}
        0 & 0 & 0 & 0 \\
0 & 1 & 0 & 0 \\
0 & 0 & -2 & 0 \\
0 & 0 & 0 & -2
      \end{pmatrix} = 
     \begin{pmatrix}
  0&0&-2&0\\
           0&0&0&0\\
           0&-1&0&-2\\
           0&1&0&4
\end{pmatrix}
\]
     
\section*{Aufgabe 4.2}
''Sind v und w Eigenvektoren von F , so ist v + w ein Eigenvektor von
F oder v + w = 0'' impliziert, dass alle Eigenvektoren zum gleichen Eigenwert gehören.  \\
\underline{Beweis durch Widerspruch.} \\
Angenommen F hat eine Basis aus Eigenwerten zu verschiedenen Eigenwerten. Folgentlich gibt es mindestens zwei linear unabhängige Vektoren zu den mindestens zwei verschiedenen Eigenwerten. Seien v und w zwei solche Vektoren. Dann gilt $v+w\neq 0$, da $v=-w$ nicht sein kann (linear unabhängig). Laut der Aufgabenstellung ist nun $v+w$ erneut ein Eigenwert von F. \\
Sei der Eigenwert zu $v$ $\lambda$ und zu $w$ $\mu$. \\
Da F linear ist, gilt 
\[F(v+w) = F(v) + F(w)\]
also
\[F(v+w) = \lambda v + \mu w\]
Da v+w ein Eigenvektor von F ist, gilt für dessen zugehörigen Eigenwert $\omega$:
\[
F(v+w) = \omega(v+w) 
\]
Also
\[
F(v+w) = \omega(v+w) = \lambda v + \mu w
\]
\[
\Rightarrow \omega = \lambda = \mu 
\]
Widerspruch, da $\lambda \neq \mu$. Folglich gibt es nur einen Eigenwert für F. \\ 
Sei dieser Eigenwert $\lambda$, dann gilt $F(v) = \lambda Id(v)$



\section*{Aufgabe 4.3}

$A^2 = \mathbb{I}_2$ bedeutet, dass die Matrix $A$ ihr eigenes Inverses ist. 
\[
\mathbb{I}_2 = \begin{pmatrix}
  1&0\\0&1
\end{pmatrix}, \quad -\mathbb{I}_2 = \begin{pmatrix}
  -1&0\\0&-1
\end{pmatrix}, \quad D^2 = \mathbb{I}_2
\]
A kann nur die Eigenwerte 1 oder -1 haben, damit $A^2 = \mathbb{I}_2$ erfüllbar ist, denn für einen Eigenvektoren v mit eigenwert $\lambda$ gilt 
\begin{align*}
    Av & = \lambda v \\
    \Rightarrow A^2 v & = \lambda^2 v \\
    Da\ A^2 = \mathbb{I}_2: \quad \mathbb{I}_2 v & = \lambda^2 v \\
    \Rightarrow v & = \lambda^2 \\
    \Rightarrow \lambda & = \pm 1
\end{align*}
Es gilt in den Fällen $A = \mathbb{I}_2$ oder $A = -\mathbb{I}_2$: 
\[
A^2 = \mathbb{I}_2 \cdot \mathbb{I}_2 = \mathbb{I}_2 
\]
\[
A^2 = -\mathbb{I}_2 \cdot (-\mathbb{I}_2) = \mathbb{I}_2 
\]
Und wenn $A = SDS^{-1}$, dann 
\[
A^2 = (SDS^{-1}) \cdot (SDS^{-1}) = SD^2S^{-1}  = S\mathbb{I}_2S^{-1} = SS^{-1} = \mathbb{I}_2
\]

\end{document}